% 後手（ごて）：自分がテンパイ寸前に誰かがリーチをかけてくる状況

\documentclass[../nankiru_main.tex]{subfiles}

% override main tex render option
\renewcommand{\renderOption}{1}

\begin{document}

\subfilepreamble{立直対応}

基準として、\textbf{子が子の立直に対してベタオリの局収支は-1100点}\citep[][p. 83]{miinin-stat}、\textbf{子が親のリーチに対してベタオリの局収支は-1700点、親が子のリーチに対してベタオリの局収支も-1700点}\citep[][p. 84]{miinin-stat}。よって、立直に対して局収支はそれらの状況でそれらの数値より高いなら全部押す。

\begin{table*}[!h]
  \centering
  \begin{tabular}{ll|SSSSSS}
     　　& 　　 & \multicolumn{2}{c}{− 8巡目 −} & \multicolumn{2}{c}{− 11巡目 −} & \multicolumn{2}{c}{− 14巡目 −} \\
    相手 & 立直の打点 & 安牌 & {10}\%牌 & 安牌 & {10}\%牌 & 安牌 & {10}\%牌 \\
    \hline
    子 & 1飜40符 & 200 & -400 & 100 & -500 & 100 & -500 \\
    　 & 2飜30符 & 700 & 100 & 500 & -100 & 500 & -100 \\
    　 & 3飜30符 & 1900 & 1200 & 1700 & 1000 & 1400 & 800 \\
    親 & 1飜40符 & -500 & \color{red}-1200 & -600 & \color{red}-1300 & -500 & \color{red}-1200 \\
    　 & 2飜30符 & 0 & -800 & -200 & -900 & -100 & -900 \\
    　 & 2飜40符 & 600 & -300 & 400 & -400 & 300 & -500 \\
    　 & 3飜30符 & 1200 & 400 & 1000 & 100 & 800 & 0
  \end{tabular}
  \caption{先制立直に対する子の追いかけ両面立直の局収支。10\%牌とは初手に10\%の危険度の牌を切った場合。先行者は2巡以上前に立直しているとする。宣言牌で振っても一発にならないとする\citep[][p. 87]{miinin-stat}}
  \label{tab:ev-ryanmen-chase}
\end{table*}

\begin{table*}[!h]
  \centering
  \begin{tabular}{lll|SSSSSS}
    　　 & 　　 & 　　　　　 & \multicolumn{2}{c}{− 8巡目 −} & \multicolumn{2}{c}{− 11巡目 −} & \multicolumn{2}{c}{− 14巡目 −} \\
    自分 & 相手 & 副露の打点 & 安牌 & {10}\%牌 & 安牌 & {10}\%牌 & 安牌 & {10}\%牌 \\
    \hline
    子 & 子 & 1飜30符 & -200 & -700 & -200 & -700 & 0 & -500 \\
    　 & 　 & 2飜30符 & 300 & -300 & 200 & -300 & 400 & -200 \\
    　 & 　 & 3飜30符 & 1200 & 500 & 110 & 400 & 110 & 400 \\
    子 & 親 & 1飜30符 & -900 & \color{red}-1600 & -900 & \color{red}-1500 & -600 & \color{red}-1300 \\
    　 & 　 & 2飜30符 & -400 & -1100 & -400 & -1100 & -200 & -1000 \\
    　 & 　 & 3飜30符 & 500 & -300 & 400 & -400 & 500 & -300 \\
    親 & 子 & 1飜30符 & -600 & -1100 & -600 & -1000 & -400 & -800 \\
    　 & 　 & 2飜30符 & -100 & -600 & -100 & -600 & 0 & -500 \\
    　 & 　 & 3飜30符 & 800 & 200 & 700 & 100 & 700 & 100
  \end{tabular}
  \caption{両面副露手で立直に対して押した時の局収支\citep[][p. 92]{miinin-stat}}
  \label{tab:ev-ryanmen-furo}
\end{table*}

\begin{table*}[!h]
  \centering
  \begin{tabular}{lll|SSSSSS}
    　　 & 　　 & 　　　　　 & \multicolumn{2}{c}{− 8巡目 −} & \multicolumn{2}{c}{− 11巡目 −} & \multicolumn{2}{c}{− 14巡目 −} \\
    自分 & 相手 & 立直の打点 & 安牌 & {10}\%牌 & 安牌 & {10}\%牌 & 安牌 & {10}\%牌 \\
    \hline
    子 & 子 & 1飜40符 & -900 & \color{red}-1300 & -800 & \color{red}-1300 & -700 & \color{red}-1200 \\
    　 & 　 & 2飜40符 & -200 & -700 & -200 & -700 & -200 & -700 \\
    　 & 　 & 3飜40符 & 600 & 0 & 600 & 0 & 400 & -100 \\
    子 & 親 & 1飜40符 & \color{red}-1800 & \color{red}-2300 & \color{red}-1600 & \color{red}-2200 & \color{red}-1400 & \color{red}-2000 \\
    　 & 　 & 2飜40符 & -1100 & \color{red}-1700 & -1000 & \color{red}-1600 & -900 & \color{red}-1500 \\
    　 & 　 & 3飜40符 & -200 & -1000 & -200 & -900 & -300 & -1000 \\
    親 & 子 & 1飜40符 & -900 & -1300 & -800 & -1300 & -700 & -1200 \\
    　 & 　 & 2飜40符 & 200 & -400 & 100 & -400 & 0 & -500 \\
    　 & 　 & 3飜40符 & 1400 & 700 & 1300 & 600 & 1000 & 300
  \end{tabular}
  \caption{愚形で立直に対して押した時の局収支\citep[][p. 98]{miinin-stat}}
  \label{tab:ev-gukei-chase}
\end{table*}

\begin{table*}[!h]
  \centering
  \begin{tabular}{lll|SSSSSS}
    　　 & 　　 & 　　　　　 & \multicolumn{2}{c}{− 8巡目 −} & \multicolumn{2}{c}{− 11巡目 −} & \multicolumn{2}{c}{− 14巡目 −} \\
    自分 & 相手 & 副露の打点 & 安牌 & {10}\%牌 & 安牌 & {10}\%牌 & 安牌 & {10}\%牌 \\
    \hline
    子 & 子 & 1飜30符 & -800 & \color{red}-1300 & -700 & \color{red}-1100 & -400 & -900 \\
    　 & 　 & 2飜30符 & -500 & \color{red}-1000 & -400 & -800 & -100 & -700 \\
    　 & 　 & 3飜30符 & 100 & -400 & 200 & -300 & 300 & -200 \\
    子 & 親 & 1飜30符 & \color{red}-1700 & \color{red}-2300 & \color{red}-1500 & \color{red}-2100 & -1100 & \color{red}-1700 \\
    　 & 　 & 2飜30符 & -1400 & \color{red}-2000 & -1200 & \color{red}-1800 & -800 & \color{red}-1500 \\
    　 & 　 & 3飜30符 & -700 & -1400 & -600 & -1300 & -400 & -1100 \\
    親 & 子 & 1飜30符 & -1300 & \color{red}-1700 & -1100 & \color{red}-1500 & -800 & -1200 \\
    　 & 　 & 2飜30符 & -900 & -1400 & -800 & -1200 & -500 & -1000 \\
    　 & 　 & 3飜30符 & -300 & -800 & -200 & -700 & 0 & -600
  \end{tabular}
  \caption{愚形副露手で立直に対して押した時の局収支\citep[][p. 101]{miinin-stat}}
  \label{tab:ev-gukei-furo}
  \vspace{2em}
\end{table*}

%--------------------------------------------
\subsection{両面聴牌}

表~\ref{tab:ev-ryanmen-chase} によると、立直の先行者が子、自分も子、そして自分の手は両面であれば\textbf{ほぼ全て押す}。自分が親の場合もほぼ全て押す\citep[][pp. 90--91]{miinin-stat}。立直の先行者は親の場合、要注意のはのみ手で危険牌を切って追いかけするとベタオリの局収支に近ついている。よって、無筋456・一発で無筋2378・無筋2378のドラを切る場合ではベタオリすべき\citep[][p. 90]{miinin-stat}。

\subsubsection*{待ちがよりいいならリスクを取れ}

\textbf{ピンフドラ三面張ぐらいあるなら10\%の放銃率があってもリスクを背負って追いかける}\citep[][pp.~121--122]{miinin-choice}。手牌\ref{mj:chase-1}は \mj{5s} 切りと \mj{8s} 切りの二択。\mj{8s} は先行者（子）の現物だが、\mj{5s} 切れば三面張はよりアガリやすい。

\tehai{23m 56788p 455678s -1m}[東1局　西家　12巡目　ドラ \mj{1m}][chase-1]

\subsubsection*{2軒立直に対する}

\textbf{ピンフのみでも追いかける}\citep[][pp.~151--152]{miinin-choice}。子の2人が立直してベタオリの局収支は$-$900、ピンフのみで追いかける局収支は$-700$、いずれにせよ損なので追いかける方がマシ。

\begin{rittai}[h]
  \centering
  
  \drawrittai{
    jicha/seat = 西,
    jicha/hand = 11123445678m 78s -5m,
    jicha/river = 6z 13p 45s 1^p,
    %
    toimen/river = 3z 9m 13s 6z 5s,
    toimen/riverb = 99p 3^z,
    %
    kami/river = 3z 9m 6z 97's 4^p,
    kami/riverb = 3^z 6^z,
    %
    shimo/hand = XXXXXXX -4'44z 77'7z,
    shimo/river = 397p 780^m,
    shimo/riverb = 1z 5^p,
    %
    kyoku = 東1局,
    kyoutaku = 1,
    dora = 7s
  }
  
  \caption{上家から立直が来て、下家がソーズの混一色の模様の場面}
  \label{rittai:ryanmen-push}
\end{rittai}

立体図~\ref{rittai:ryanmen-push} は両軒立直に近い場況、\mj{4m} を切ったら \mj{69s} 待ちの立直・ドラ1になる。

まずは一発の巡目に \mj{0m} 切りを敢行する下家に \mj{78s} を切ることはない。このままで \mj{4m} を切って立直していいが、安全牌の \mj{5m} をツモ切って、ツモ \mj{9m} による一通やツモ \mj{7s} で危険な \mj{4m8s} で勝負することもあり\citep[][pp.~106--109]{doi}。

\subsubsection*{ダマ・副露}

全ての状況で両面ダマの局収支は両面立直より著しく低い\citep[][p. 92]{miinin-stat}。役があってもがんがんと立直をかけよう。

両面副露手で立直に対してがんがんと押すのはどうでしょう。表~\ref{tab:ev-ryanmen-furo} によると、多くの場合では押せる。ベタオリすべきな場合は子で親の立直に対して1飜だけの副露手。

\subsubsection*{1000点の押し引き}

\begin{rittai}[h]
  \centering
  
  \drawrittai{
    jicha/seat = 北,
    jicha/hand = 22345m 56p 567s -4'35m,
    jicha/river = 1m 51z 81s 3z,
    jicha/riverb = 3z 7m 9p 8m 3s,
    %
    toimen/river = 2z 1s 4^z 5m 38^p,
    toimen/riverb = 1^m 1^z 6^s 1^z 2^p 2^z,
    %
    kami/river = 1z 1p 8m 235z,
    kami/riverb = 66z 8p 11m,
    %
    shimo/river = 3z 9s 12p 6z 8^m,
    shimo/riverb = 11s 2m 6s 3p 6z,
    %
    kyoutaku = 1,
    dora = 7m
  }
  
  \caption{1000点の聴牌はどこまで押すか}
  \label{rittai:cheap-tenpai}
\end{rittai}

立体図~\ref{rittai:cheap-tenpai}の通り、自分が1000点の聴牌。通ってない筋は \mj{3m}--\mj{6m}--\mj{9m}、\mj{1p}--\mj{4p}--\mj{7p}、\mj{2s}--\mj{5s}--\mj{8s} と \mj{4s}--\mj{7s}。ツモ (1) \mj{3m} (2) \mj{2s} (3) \mj{5s} (4) \mj{9m} ならどうするか、という課題。対面の河による通っていない筋を分析すると\citep[][pp.~128--133]{tetsuoshi}、

\begin{itemize}
\item \mj{36m}：立直前に切られているので他の筋より少しマシ
\item \mj{47s}：\mj{6s} が3枚見えるので他の筋より少しマシ\footnote{\citet[p.~132]{tetsuoshi}によると \mjfn{6s} も立直前に切られているが、実際にその \mjfn{6s} は南家の河に写っていない。}
\item \mj{9m}：\mj{8m} が3枚見えるのでワンチャンス。けれど他の2人がベタオリしているのにラストの \mj{8m} が出てこないので南家の手にある可能性が高い
\item \mj{147p25s}：特に危険
\end{itemize}

よって、のみ手なら \mj{3m} まで押す。3900点あるなら \mj{9m} で押すか微妙だが、8000点あるなら全て押す。

%--------------------------------------------
\subsection{愚形聴牌}

表~\ref{tab:ev-gukei-chase} によると、親の立直が来て自分の手は愚形リーのみならオリる。立直を含めて2飜があっても切る牌が危険ならオリてもいい。自分が親でのみ手から切る牌が無筋456ならオリる。その他の場合では押す。

愚形副露手なら状況が厳しくなる。表~\ref{tab:ev-gukei-furo} によると自分が子でも親でも立直に対して危険牌で1飜の手は押すべきではない。親の立直に対して2飜があっても切る牌が危険ならオリる。

ダマの場合は立直の場合より大差が出ないのでダマより立直で打点を高くする方がいい\citep[][p. 103]{miinin-stat}。


\subsubsection*{七対子}

\textbf{親の立直に対しても1枚切れ字牌単騎ぐらいなら押す}\citep[][pp.~133--134]{miinin-choice}。七対子のみでも追いかける\footnote{待ちがそんなによくないというデータはまだないのでこの辺はまだわからない。}。


\subsubsection*{回し打ち}

\textbf{1300のみ手で危険牌を切らないと聴牌とれないなら聴牌を崩す}\citep[][pp.~137--138]{miinin-choice}。例え\ref{mj:chase-2}のように \mj{9m} が現物なら \mj{9m} を切って良形聴牌を目指す。特に連続形があれば回し打ちは有効。

\tehai{1179m 222567p 3456s}[東1局　西家　12巡目　ドラ \mj{1s}][chase-2]

2600のみ手ならどちらでもいい。場況・河・点数によって決める。


\subsubsection*{2軒立直に対する}

\textbf{ドラが2枚あってもオリる}\citep[][pp.~153--154]{miinin-choice}。5200で2軒立直と勝負する局収支は$-1000$しかないから。


\subsubsection*{データの補足} 

ここまでの話は安牌と10\%牌に基づく。空間が少ないため全ての場合を挙げるのは不可能。ですが、内挿（ないそう）・外挿（がいそう）によると5\%牌と15\%牌の局収支を補足できる。例え、表~\ref{tab:ev-ryanmen-chase}を参照して、手牌\ref{mj:3-stat-1}から \mj{4p} （15\%牌）を切って立直する局収支は$-1300 + (-1300 - (-600)) / 2 = -1650$。ベタオリの局収支（$-1700$）と大差がないので、場況（点数状況・待ち枚数など）によって判断する\citep[][pp. 93--95]{miinin-stat}。

\tehai{123666m 478p 23466s}[東1局　南家　親が立直　11巡目　ドラ \mj{9m}][3-stat-1]

切る牌がドラ、または一発の場合、局収支は元の数値から宣言牌の危険度に放銃平均値の差をかけた値を引く。相手が子の場合、放銃平均値の差は2300点、親なら3400点。よって、親の立直に対してドラの10\%牌を切ったら局収支が340点低くなる\citep[][p. 95]{miinin-stat}。


\subsubsection*{残り2枚しかない愚形}

節~\ref{subsec:senseigukei} で先に触れているが、ここで実際の場面を紹介にいく。便宜上に残り2枚しかない愚形の立直判断はここで再掲する。

\begin{table*}[t]
  \centering
  \begin{tabular}{l|lll}
    ドラ枚数 & いい形がない & いい形がある & 先制された場合 \\
    \hline
    なし & ダマ & 聴牌外し & 追いかけない \\
    1枚 & 待ちがいいなら立直する & 待ちがいいなら立直する & 追いかけない \\
    2枚 & 立直 & 待ちがいいなら立直する & 待ちがいいなら立直する \\
    3枚以上 & 立直 & 立直 & 立直
  \end{tabular}
  \caption{愚形待ちが残り2枚しかない状況の立直判断（七対子除き）}
\end{table*}

立体図~\ref{rittai:twoleft} はMリーグのある局面。自分が親でドラ2愚形聴牌だが待ちの \mj{5p} が既に2枚切れ。しかも対面から立直来ている。残りの \mj{5p} が山にありそうか否か分かりにくいが、上がればほぼ満貫なので、ここで対面と勝負する。

\begin{rittai}[h]
  \centering
  
  \drawrittai{
    jicha/seat = 東,
    jicha/point = 28500,
    jicha/hand = 123m 12346p 45678s -0s,
    jicha/river = 9m 1s 3z 9p 8^p 2^z,
    jicha/riverb = 4z 7z 7^z 6^z 2^s 1z,
    %
    toimen/point = 18600,
    toimen/river = 5z 7z 5^z 3^z 4s 2^z,
    toimen/riverb = 5s 6p 7^s 8m 6^m 9'm,
    %
    kami/point = 48400,
    kami/river = 2z 6m 6z 5^p 3^z 3m,
    kami/riverb = 3^5^6^1^p 9^s 5^z,
    %
    shimo/point = 3500,
    shimo/river = 1z 8m 6s 2^z 5m 7^s,
    shimo/riverb = 2s 6^s 9^m 8p 4p 6^m,
    %
    kyoku = 南3局,
    kyoutaku = 1,
    dora = 2m
  }
  
  \caption{親で2枚切の嵌 \mj{5p} で聴牌の場面}
  \label{rittai:twoleft}
\end{rittai}


%--------------------------------------------
\subsection{現物待ち}

\textbf{両面・3飜以上}、または\textbf{愚形}の場合ではダマの方が有利\citep[][pp. 139--142]{miinin-choice}。表~\ref{tab:ev-genbutsu} によると、両面の場合、手牌の価値が高いほど最善手がダマに寄るが、他の要素（点数状況、残り枚数など）による判断も重要。なお、愚形の場合ならほとんどダマの方が正解。

\begin{table}[ht]
  \centering
  \begin{tabular}{llSS}
    待ち & 手牌基本価値 & 立直 & ダマ \\
    \hline
    両面 & 1飜30符 & 800 & 600 \\
    両面 & 3飜30符 & 3400 & 3600 \\
    愚形 & 1飜40符 & -100 & 100 \\
    愚形 & 3飜40符 & 1000 & 1800 \\
  \end{tabular}
  \caption{現物待ちで現物を切って聴牌を取る11巡目の局収支}
  \label{tab:ev-genbutsu}
\end{table}


%--------------------------------------------
\subsection{序盤一向聴}

序盤では通っている筋の本数はまだ少ないので、押し引きはちょっと押し寄りであろう。切った牌が少ないため、安全度の読みはあまりできない。手牌\ref{mj:3-1-1}から\ref{mj:3-1-4}までは全て立体図 \ref{rittai:3-1} の状況を仮定している\citep[][pp. 16--21]{tetsuoshi}。

\begin{rittai}[h]
  \centering
  
  \drawrittai{
    jicha/seat = 北,
    jicha/hand = ??????????????,
    jicha/river = 1z 1m 9s 8p 3z,
    %
    toimen/river = 2z 9s 4z 2m 3'p 9^p,
    %
    kami/river = 1p 1z 2z 8s 5z 9p,
    %
    shimo/river = 3z 1m 5z 2m 7p 3z,
    %
    dora = 6p
  }
  
  \caption{序盤の立体図}
  \label{rittai:3-1}
\end{rittai}

\tehai{340788m 23789p 11 -7s}[][3-1-1]

\textbf{ドラ1の完全一向聴は押す}。手牌\ref{mj:3-1-1}の価値が3900ある上受け入れが最大限なので押す\citep[][pp. 16--21]{tetsuoshi}。ピンフドラ1があっても押さない状況は、

\begin{itemize}
\item 親リー
\item ピンフのみ
\end{itemize}

\tehai{3407m 123789p 117 -3s}[][3-1-2]

くっつき一向聴ならどちらでもいい。手牌\ref{mj:3-1-2}で \mj{3s} を押した後カン \mj{8m} になったら待ちがそこそこだが、中張の \mj{6m6s} や1枚切れの \mj{8s} は流石に微妙。嵌張で追いかけるかどうかはそのときの場況・他家の動向で判断してもいい。また、良形が不確定ば分、\mj{1s} でオリるのもある。\citep[][pp. 22--25]{tetsuoshi}。

\tehai{307m 123789p 1167 -3s}[][3-1-3]

\textbf{ドラ1けど愚形＋良形ならオリる}。ドラ2ならまだ押せる。だとしても、手牌\ref{mj:3-1-3}の手牌なら中筋のカン \mj{4m} で待つのも悪くない。対面が \mj{2m} 切っているので \mj{1m} は出やすいし、\mj{7m} が通れば \mj{4m} も出やすい\footnote{この推理は残りの二人が中筋を切る意向があるという前提で成り立つ。}\citep[][pp. 26--29]{tetsuoshi}。

\tehai{34m 123789p 56999 -3s}[][3-1-4]

\textbf{受け入れの広さに関わらずのみ手はオリる}。手牌\ref{mj:3-1-4}の受け入れは打 \mj{3s} で8種の \mj{2345m4567s} があるが、価値の最低限は1300、裏ドラ乗っても打点の上昇幅はそんなに大きくないので、オリる方は若干有利\citep[][pp. 30--35]{tetsuoshi}。

\subsubsection*{副露手}

\tehai{222340m 56p 3667s -4'35m}[][3-2-1]

同じの立体図 \ref{rittai:3-1} を想定して、今回は\ref{mj:3-2-1}のような3900一向聴の副露手。副露手は一発・裏ドラ・ツモがつかないため平均打点は門前手より低いだが、鳴けるので門前手より早い。打点が2000としても序盤で残り筋が多いなら押す\citep[][pp. 46--49]{tetsuoshi}。

\subsubsection*{七対子}

\begin{rittai}[h]
  \centering
  
  \drawrittai{
    jicha/seat = 北,
    jicha/hand = 22057m 33p 11899s 7z -6p,
    jicha/river = 1z 1m 2s 8p 3z,
    %
    toimen/river = 2z 1s 4z 2m 3'p 4^m,
    %
    kami/river = 1p 1z 2z 8s 5z 9p,
    %
    shimo/river = 3z 1m 5z 2m 7p 3z,
    %
    kyoutaku = 1,
    dora = 5p
  }
  
  \caption{序盤で自分の手牌が七対子一向聴の場合}
  \label{rittai:chiitoi}
\end{rittai}

七対子一向聴ならオリ寄り。立体図\ref{rittai:chiitoi}であれば、ドラが重なったり、待ち牌が場況によっていい待ちであったり、序盤で切る牌が嵌張（\mj{68m}）・辺張（\mj{89m}）に放銃するリスクを負って押すのが少し有利。ただし、これらの条件が整わないなら \mj{1s} でオリよう\citep[][pp. 50--55]{tetsuoshi}。

\subsubsection*{安め高めがある手牌}

\begin{rittai}[h]
  \centering
  
  \drawrittai{
    jicha/seat = 北,
    jicha/hand = 222478m 6789p 367 -8s,
    jicha/river = 1z 1s 2s 8p 3z,
    %
    toimen/river = 2z 2s 4z 2m 3'p 4^m,
    %
    kami/river = 1z 2z 8s 5z 5z 4z,
    %
    shimo/river = 3z 1p 5z 7p 6z 3p,
    %
    kyoutaku = 1,
    dora = 3z
  }
  
  \caption{序盤で手牌の価値の範囲が大きい場合}
  \label{rittai:3-2}
\end{rittai}

立体図 \ref{rittai:3-2} を想定する。自分の手は高めで立直・タンヤオ・ピンフ・三色の8000（一発・裏ドラ・ツモのどちらがあったら12000）、安めで立直のみの1300。ここで現物の \mj{4m} を切ると \mj{3m} の受け入れと \mj{5m} の入れ替えが無くなり、打点が一気に下がるので、序盤で \mj{3s} 切りは有利。以降の巡目で \mj{6m} を引いたらダマにしよう\citep[][pp. 62--65]{tetsuoshi}。

\subsubsection*{切る牌が本命牌}

\tehai*{7z 8s 2p 5m 3^z 7'm 1^z}[対面（南家）の河][]

\tehai{22566788p 34678s -9m}[自分の手牌(1)　ドラ \mj{4p}][3-6-1]

\tehai{22566788p 34789s -9m}[自分の手牌(2)　ドラ \mj{4p}][3-6-2]

\mj{5m} を先に切って \mj{7m} で立直を宣言するのは、もともと手牌の中に \mj{5799m} や \mj{5778m} の形がある確率が非常に高い。そして自分の手牌は、\ref{mj:3-6-1}なら最大限で立直・タンヤオ・一盃口・ピンフで7700以上（ツモれば8000以上）、\ref{mj:3-6-2}ならタンヤオがつかないためほぼ3900である。ですので、\ref{mj:3-6-1}を押すのはちょっと微妙であるが、\ref{mj:3-6-2}を押すのは流石に無理である\citep[][pp.~66--73]{tetsuoshi}。

\subsubsection*{切る牌が生牌ドラ字牌}

天鳳位による討論でも意見が結構割れていても、\textbf{序盤で満貫良形一向聴なら切ってもいい}\citep[][pp.~114--115]{tetsuoshi}。

\subsubsection*{安牌がない時}

\tehai*{5z 829^p 3^z 2s 6'p 9^m}[対面（南家）の河][]

\tehai{667m 345p 1234599s -3m}[自分の手牌　ドラ \mj{4p}][3-7-1]

自分の手は3900の完全一向聴だが、現物は \mj{2s} しかない。手牌\ref{mj:3-7-1}で \mj{3m} をツモ切りすると後 \mj{6m} を切って立直すると \mj{6m} が中筋になるのでいい押しと考えられる。しかしツモった牌は \mj{7s} などのような他の無筋ならリスクが高すぎるので \mj{2s} でオリる方がいい。相手が親、または自分の手がピンフのみならオリる方がいい\citep[][pp.122--127]{tetsuoshi}。

\subsubsection*{立直後のポン}

\begin{rittai}[h]
  \centering
  
  \drawrittai{
    jicha/seat = 東,
    jicha/hand = 122056m 05678p 56s,
    jicha/river = 6z 2p 4^z 9s,
    jicha/point = 26600,
    %
    toimen/river = 3s 7^z 8^4^p,
    toimen/point = 35800,
    %
    kami/river = 1p 9s 6^z 5'p,
    kami/point = 14200,
    %
    shimo/river = 43z 1s 5z,
    shimo/point = 22400,
    %
    kyoku = 南1局,
    honba = 5,
    kyoutaku = 1,
    dora = 6z
  }
  
  \caption{立直を受けている状態でポンするか}
  \label{rittai:riipon}
\end{rittai}

立体図\ref{rittai:riipon}のように、立直の先行者が暫定最後位にも関わらず\textbf{親で好形高打点一向聴ならSuphxは \mj{5p} をポンして \mj{1m} で押すとした}\citep[][pp.~46--47]{suphx}。立直を受けている状態で、ポンするという行為自体が安牌2枚を消費するので、以下の条件がいずれも満たさないならまず見受けられない。

\begin{itemize}
\item 聴牌
\item 安牌多数持ち
\item 好形高打点一向聴
\end{itemize}

同卓のメンツのレベルが上がるほど聴牌か高打点かのパターンが絞られる。そういうとき押せる一向聴の手牌も押せなくなる。


%--------------------------------------------
\subsection{中盤一向聴}

\subsubsection*{完全一向聴}

\begin{rittai}[h]
  \centering
  
  \drawrittai{
    jicha/seat = 北,
    jicha/hand = 340788m 23789p 11 -3s,
    jicha/river = 1z 1m 9s 8p 3z 7^s,
    jicha/riverb = 1^m 3^p 7^z 7^z 4^z,
    %
    toimen/river = 2z 9s 4z 2m 3'p 9^p,
    toimen/riverb = 1^z 6^p 6^m 1^p 5^s 8^m,
    %
    kami/river = 1p 1z 2z 8s 5z 9p,
    kami/riverb = 2m 6p 3z 4s 5s 4z,
    %
    shimo/river = 3z 1m 5z 2m 7p 3z,
    shimo/riverb = 4s 3p 2^z 9m 2s 2s,
    %
    kyoutaku = 1,
    dora = 5p
  }
  
  \caption{中盤の立体図}
  \label{rittai:chuuban}
\end{rittai}

立体図 \ref{rittai:chuuban} は立体図 \ref{rittai:3-1} に続いて、6回のツモ切りもしてもまだ一向聴の状態。残り筋が6本（ \mj{47m}、\mj{25p}、\mj{47p}、\mj{58p}、\mj{36s}、\mj{69s} ）で、\mj{3s} を切ったら$1/6\approx17.8\%$の確率で放銃する。巡目はもう中盤〜終盤なので和了率はそこまで高くないので一旦 \mj{8m} を切ってオリよう。もし \mj{2s} を引いたら \mj{14p} の闇聴にしよう\citep[][pp. 40--45]{tetsuoshi}。

\begin{table}[ht]
  \centering
  \begin{tabular}{lllllll}
    筋    & 14 & 25 & 36 & 47 & 58 & 69 \\
    \hline
    マンズ & $\times$ & $\times$ & $\times$ & & $\times$ & $\times$ \\
    ピンズ & $\times$ & & $\times$ & & & $\times$ \\
    ソウズ & $\times$ & $\times$ & & $\times$ & $\times$ &
  \end{tabular}
  \caption{全ての筋のなかでまだ6本の筋が通っていない。$\times$マークがついているのは通っている筋}
\end{table}

\begin{rittai}[h]
  \centering
  
  \drawrittai{
    jicha/seat = 西,
    jicha/hand = 55m 788p 23456789s -0p,
    jicha/river = 4z 1m 2p 3^z 1^p 8m,
    jicha/riverb = 8m 5^z 4p 9^s,
    %
    toimen/river = 43z 289p 4m,
    toimen/riverb = 6^p 6s 3^p 77z,
    %
    kami/river = 52z 12m 64p,
    kami/riverb = 6s 6m 1^p 3s 1^z,
    %
    shimo/river = 31z 9s 2p 7m 3s,
    shimo/riverb = 1s 8^s 3'm 5^z,
    %
    kyoutaku = 1,
    dora = 1s
  }
  
  \caption{超良形一向聴で片筋の \mj{0p} をツモった中盤}
  \label{rittai:good-wait}
\end{rittai}

立体図~\ref{rittai:good-wait} では一通も見込む超良形一向聴だがここで \mj{0p} を掴まえてしまった。まずはどれぐらい良形でも片筋 \mj{0p} を放つことはない。一旦 \mj{8p} を切って \mj{58p} の筋を試す選択肢もあるが、もう11巡目に加えて \mj{6p} が既に2枚切れなので結局オリになるケースが多い。点棒が切実に必要ではないなら \mj{89s} でオリる\citep[][pp.~126--129]{doi}。



\subsubsection*{両面＋両面一向聴}

\begin{rittai}[h]
  \centering
  
  \drawrittai{
    jicha/seat = 西,
    jicha/hand = 2344m 2367p 12345s -6s,
    jicha/river = 53z 98m 7s 3^z,
    %
    toimen/hand = XXXXXXXXXX -4'44p,
    toimen/river = 1m 8p 5z 6p 4m 6m,
    toimen/riverb = 8^s,
    %
    kami/river = 4z 18m 57p 6's,
    kami/riverb = 6^m,
    %
    shimo/river = 67z 2p 7^m 7p,
    %
    kyoutaku = 1,
    dora = 3s
  }
  
  \caption{メンピンドラ1までどう押すか}
  \label{rittai:chuuban2}
\end{rittai}

立体図~\ref{rittai:chuuban2} で自分はピンフドラ1の一向聴。\mj{3m} は筋、\mj{7p} は現物、\mj{236p} はともにワンチャンス、そして自分の手牌は今搭子オーバーフローなので、打 \mj{23m} と打 \mj{67p} の二択。

\mj{23m} を落とすのは \mj{67p} より危険だが、\mj{14m} もう薄くなっている上、ピンズの場況がいいので、\mj{23m} 落としていい。まだ、安全度重視なら \mj{7p} から切って、聴牌したら場況によりリーチかダマか考えるのもある。ドラが1枚増えたら \mj{23m} 切りに寄る。

ちなみに親の捨て牌に「ワタシトイトイヤッテイマス」と書いてある。ですが自分がピンフドラ1一向聴なので一旦無視しよう\citep[][pp.~110--113]{doi}。

\subsubsection*{両嵌＋両面一向聴}

\begin{rittai}[h]
  \centering
  
  \drawrittai{
    jicha/seat = 西,
    jicha/hand = 135m 123789p 2399s -6p,
    jicha/river = 47z 8m 2p 6s 3^z,
    jicha/riverb = 8s 4^z 1m,
    %
    toimen/river = 1s 8p 75z 4p 3^s,
    toimen/riverb = 974m 1^p,
    %
    kami/river = 3z 1p 831m 5p,
    kami/riverb = 6^z 3^m 42s,
    %
    shimo/river = 26z 3s 24m 7s,
    shimo/riverb = 4^s 4'p 1^z,
    %
    kyoutaku = 1,
    dora = 5m
  }
  
  \caption{高目純チャン三色だが無筋を掴まえてしまった局面}
  \label{rittai:takame-iishanten}
\end{rittai}

立体図~\ref{rittai:takame-iishanten} で自分の手牌は両嵌張＋両面で受け入れが16枚しかないだが高目が純チャンか、三色か、もしくはピンフかの勝負手。

ですがこんな手でも下家の河にドラ \mj{5m} や無筋 \mj{6p} を放つ価値がない。\mj{37s} などの中張牌をいっぱい切ったのは手牌が整った有力な証拠。そしてドラ \mj{5m} も1枚しか見えないので下家の手牌に複数枚あっても珍しくない。ここで一旦 \mj{1m} を切って、\mj{4m} が入ったら \mj{9p} で勝負する。ただし以外の危険牌を掴まえたらオリる\citep[][pp.~122--125]{doi}。


\subsubsection*{格が違う向聴戻し}

\begin{rittai}[h]
  \centering
  
  \drawrittai{
    jicha/seat = 東,
    jicha/hand = 67m 66p 334445678s -9m,
    jicha/river = 2m 89s 4^m 2^z 2^m,
    jicha/riverb = 24p 1^s 9^p 6s,
    jicha/point = 25000,
    %
    toimen/river = 431z 9s 8^p 5s,
    toimen/riverb = 1m 2^z 6^p 6'z,
    toimen/point = 25000,
    %
    kami/hand = XXXXXXXXXX -2'34s,
    kami/river = 5z 2^1z 9s 9p 3z,
    kami/riverb = 7^p 738^1m,
    kami/point = 25000,
    %
    shimo/river = 5z 8^m 43z 2s 4^z,
    shimo/riverb = 1^s 9p 2z 81^s,
    shimo/point = 24000,
    %
    kyoku = 東3局,
    kyoutaku = 1,
    dora = 6z
  }
  
  \caption{中盤で無筋を掴まえて向聴数を戻したい局面}
  \label{rittai:shanten-back}
\end{rittai}

\textbf{複合している部分を残すのは基本}\citep[][pp.~48--49]{suphx}。立体図~\ref{rittai:shanten-back} の場合、ドラ \mj{6z} 立直に対して無筋の \mj{9m} が掴まえて一発を避けるため現物の \mj{6p} と \mj{5s} を切る二択。Suphxは打 \mj{6p} とした。そうするとツモ \mj{4578p} の変化が拾えるし、切りにくい \mj{17s} も拾える。一方、\mj{5s} にすれば聴牌しても残り2枚の双碰待ち、ケイテンではないなら微妙。

\subsubsection*{マシな無筋}

\begin{rittai}[h]
  \centering
  
  \drawrittai{
    jicha/seat = 西,
    jicha/hand = 567m 22399p 07s 666z -6m,
    jicha/river = 2z 8p 2^3^m,
    jicha/point = 8200,
    %
    toimen/river = 42z 18p 51z,
    toimen/riverb = 9'p,
    toimen/point = 42100,
    %
    kami/river = 175z 771m,
    kami/riverb = 1's,
    kami/point = 8500,
    %
    shimo/hand = XXXXXXX -0'34p3'24s,
    shimo/river = 1^3z 1p 5z 2p 3^z,
    shimo/point = 39200,
    %
    kyoku = 南2局,
    honba = 1,
    kyoutaku = 2,
    dora = 4p
  }
  
  \caption{無筋を押さなければならない局面（北家は \mj{3s} チー打 \mj{1p}、\mj{0p} チー打 \mj{5z}）}
  \label{rittai:better-musuji}
\end{rittai}

立体図~\ref{rittai:better-musuji} のように、自分が暫定最後位、そして三位の南家から立直が入ってきて、北家も副露してドラを2枚晒している状態、\textbf{Suphxは北家の現物の \mj{2p} を切った}\citep[][pp.~50--51]{suphx}。ゼンツにすればこの \mj{6m} はツモ切りだろう。オリるなら \mj{6z} の3連打もあり。\mj{2p} は無筋で「中途半端」の押しに見える。

ここの最大な不確定要素は北家。\mj{0p} チー後で打 \mj{2p} のは食い替えたばかり、速度はわかりにくいところ。よって、ここは一旦 \mj{2p} を切って一向聴を維持しつつ、次の引く牌や北家の動向次第で行動する。ちなみにこの直後Suphxは \mj{1p} チー打 \mj{6m} とした。

\subsubsection*{手牌維持}

\begin{rittai}[h]
  \centering
  
  \drawrittai{
    jicha/seat = 東,
    jicha/hand = 340799m 78s 55z -4p -5'05s,
    jicha/river = 3z 9p 1z 2s 1^2^p,
    jicha/riverb = 9^6p,
    jicha/point = 18400,
    %
    toimen/river = 9m 4^6z 2s 7z 3^s,
    toimen/riverb = 9s 1^p,
    toimen/point = 36400,
    %
    kami/river = 4z 8^m 3z 2^3^m 4^z,
    kami/riverb = 2'z,
    kami/point = 36300,
    %
    shimo/river = 1^p 9s 6z 2s 1^m 3s,
    shimo/riverb = 3^s 1^z,
    shimo/point = 7900,
    %
    kyoku = 南2局,
    kyoutaku = 1,
    dora = 2p
  }
  
  \caption{親一向聴だが立直がくる場面}
  \label{rittai:oya-maintain}
\end{rittai}

立直対して親のSuphxは外側の \mj{9m} を切った\citep[][pp.~52--53]{suphx}。親番が終われば残り2局、そして一位との点差が大きいのでこの点棒状況ならオリによるだろう、もしラス目の下家を飛ばせてくれば何卒。けれどここは5800の一向聴、アガればライス回避はほぼ確実。

アガりに向かうけれどSuphxはド無筋の \mj{7m} や \mj{4p} でゼンツはしない。やや安全そうな \mj{9m} で押し、一向聴を維持する。聴牌なら押し、安牌が増えたらオリ、頭ができないならオリる。とはいえ、手牌がもっと悪い場合ちゃんと \mj{3m} でオリよう。

\begin{rittai}[h]
  \centering
  
  \drawrittai{
    jicha/seat = 西,
    jicha/hand = 246678m 44p 34567s -4p,
    jicha/river = 45z 97p 1^s 7^z,
    %
    toimen/hand = XXXXXXXXXX -11'1z,
    toimen/river = 37z 86s 4^z 3p,
    toimen/riverb = 4^m 2p,
    %
    kami/river = 9p 1^m 89s 6p 5m,
    kami/riverb = 3'z 7^m,
    %
    shimo/river = 3z 742m 6p 9s,
    %
    kyoku = 東1局,
    kyoutaku = 1,
    dora = 2s
  }
  
  \caption{立直とダブ東が入った場面}
  \label{rittai:maintain}
\end{rittai}

立体図~\ref{rittai:maintain} は自分がタンヤオ一向聴だが、上家が立直してきて、対面がダブ東を晒している場面。東1局の平場だが、ドラや赤ドラは一切見えない。

オリに寄るなら打 \mj{4m}。\mj{4m} は親の現物だし、上家の立直宣言牌は字牌なので中筋が刺されることほとんどない。安全の一向聴を維持するため打 \mj{7m} もあるが、打 \mj{4m} の受けがより広い。

押しに寄るなら打 \mj{2m}。\mj{2m} は上家の筋だし、数巡前に下家の \mj{2m} が鳴かれなかったら \mj{2m} が対面にも通しそう。三面張といういい形がある上、ドラ受けもあるので、聴牌に至るとカン \mj{5m} でも押す\citep[][pp.~102--105]{doi}。

\subsubsection*{吸収性能}

\begin{rittai}[h]
  \centering
  
  \drawrittai{
    jicha/seat = 南,
    jicha/hand = 340557m 456p 367s 777z,
    jicha/river = 9p 1s 1p 6z 9^p 2p,
    jicha/riverb = 2z 5p,
    jicha/point = 26000,
    %
    toimen/river = 23^74z 9m 7p,
    toimen/riverb = 8p 1m,
    toimen/point = 24500,
    %
    kami/river = 2z 9p 51z 13m,
    kami/riverb = 6m 5'p 6m,
    kami/point = 19600,
    %
    shimo/river = 2z 9m 1^z 9p 6z 9s,
    shimo/riverb = 5z 1m,
    shimo/point = 28900,
    %
    kyoku = 東3局,
    kyoutaku = 1,
    dora = 4p
  }
  
  \caption{親の立直を受けている場面}
  \label{rittai:absorb}
\end{rittai}

\textbf{同じ無筋でもSuphxは打 \mj{7m}}\citep[][pp.~54--55]{suphx}。立直者の河に切られている色に偏りがある際、その逆色は特に危険。今回立体図~\ref{rittai:absorb} の場合、\mj{7m} は立直の前に固定された \mj{67m} にしか刺さらないのに対し、\mj{3s} は両面・辺張・嵌張・双碰いずれも否定されていない。ここで打 \mj{7m} としたらソーズを引いたら \mj{34m} を払ってまだ聴牌の可能性が残る。


%--------------------------------------------
\subsection{両向聴以上}

\subsubsection*{壁か暗刻落としか}

\tehai{67m 40p 224468s 444z -1p}[東1局　東家　ドラ \mj{6p}　\mj{2p} が4枚見えている]

8巡目で下家が立直してきて、\mj{2p} が4枚見えているので \mj{1p} がかなり安全だが\textbf{Suphxは \mj{4z} の暗刻を落とした}\citep[][pp.~44--45]{suphx}。現物は \mj{4z} のみ。\mj{1p} で押すとしたら無筋 \mj{2468s} が出ていく可能性が高いので、暗刻落としによって \mj{2468s} を浮かせてよりいい聴牌を目指す。暗刻落としのうちに安牌が増えてきて \mj{1p} すら最後まで切らないことも可能。

\end{document}